\chapter{Interpretations}
\label{chap:interpretations}

% ── Introduction ──────────────────────────────────────────────────────────────
\section{Introduction}
\label{sec:interp-intro}

The mathematical formalism developed in the preceding chapters admits several
physical interpretations.  In this chapter we examine the most prominent
interpretations of quantum torsion in the context of the Quantum Observable
Universe and discuss their philosophical and physical ramifications.

% ── Copenhagen Interpretation ─────────────────────────────────────────────────
\section{Copenhagen-Style Interpretation}
\label{sec:interp-copenhagen}

% TODO: articulate how the QOU torsion manifests under a Copenhagen reading
%       of quantum measurement and collapse.

\begin{remark}
  In the Copenhagen picture, the quantum torsion $\mathcal{T}^\nabla$
  encodes the non-commutativity of successive measurements.  Each
  measurement event introduces a torsional contribution to the effective
  connection on the observable manifold.
\end{remark}

% ── Many-Worlds Interpretation ────────────────────────────────────────────────
\section{Many-Worlds Interpretation}
\label{sec:interp-mwi}

% TODO: discuss how the torsion lifting functor $\mathcal{L}$ relates to
%       branching in the Everettian picture.

% ── Relational QM ─────────────────────────────────────────────────────────────
\section{Relational Quantum Mechanics}
\label{sec:interp-relational}

% TODO: formulate the QOU torsion in Rovelli's relational framework,
%       where torsion captures relational differences between observer frames.

% ── Categorical / Topos Interpretation ───────────────────────────────────────
\section{Categorical and Topos-Theoretic Interpretation}
\label{sec:interp-topos}

% TODO: interpret the higher-categorical lifting of Chapter~\ref{chap:torsion}
%       within a topos-theoretic model of quantum logic.

\begin{definition}[Quantum Topos]
  A \emph{quantum topos} is a Grothendieck topos $\mathcal{T}$ equipped with
  an internal quantum logic given by a Heyting algebra of propositions derived
  from the spectral sheaf of $\mathcal{A}$.
\end{definition}

% ── Comparison ────────────────────────────────────────────────────────────────
\section{Comparison of Interpretations}
\label{sec:interp-comparison}

% TODO: tabulate the predictions and mathematical differences among the
%       interpretations surveyed above.

% ── Summary ───────────────────────────────────────────────────────────────────
\section{Summary}
\label{sec:interp-summary}

% TODO: summarize the interpretational landscape and motivate the experimental
%       tests presented in Chapter~\ref{chap:experimental}.
